% Limitations
% viscode_shared_subspace_probe — D028 framing
% Generated: 2026-04-11

\section{Limitations}
\label{sec:limitations}

\paragraph{(a) Sample size and model diversity.}
The headline Qwen analysis operates on $N{=}12$ cells, closer to a within-family ablation than a cross-model comparison.
The cross-architecture extension (\S\ref{sec:generalizability}) adds three models from distinct families (Mistral, BigCode, MoE), confirming that the weak signal generalizes.
However, all models are sampled at the same $n{=}252$ triples with 7 equidistant layers, which may compress inter-model differences.
Additional scales (e.g., 70B+), training paradigms (e.g., RLHF-heavy models), and non-Western tokenizers remain untested.

\paragraph{(b) Power interpretation.}
A1 power = 1.0 confirms that the permutation test can detect a CKA signal of the observed magnitude; it does \emph{not} establish that the effect size is large enough to be practically meaningful or to support claims about visual semantics.
The distinction between statistical significance and practical significance is critical given the small absolute CKA values (0.08--0.15).

\paragraph{(c) Contamination.}
VisCoder2 is fine-tuned on VisCode-Multi-679K, which includes VCM-Asy (overlapping with the probe pool source for Asymptote) and VCM-TikZ (71.7\% overlap with the TikZ probe pool source).
This creates a memorization floor for VisCoder2's Asy- and TikZ-related cells.
The headline $N{=}12$ statistic excludes VisCoder2 from the primary comparison as mitigation.
However, the remaining Qwen2.5 models may also have been exposed to SVG and TikZ content from Common Crawl during pretraining, a risk that cannot be fully quantified.

\paragraph{(d) No visual-level pairing.}
The probe pool is semantically matched via SBERT caption similarity (cosine $\geq 0.70$), not by rendering the same visual content in multiple formats.
True visual-level paired datasets---where identical geometric scenes are expressed in SVG, TikZ, and Asymptote---do not currently exist at sufficient scale.
Without such data, we cannot fully disentangle visual semantic similarity from textual and syntactic similarity in the CKA measurement.

\paragraph{(e) Scope of formats.}
We examine three visual code formats (SVG, TikZ, Asymptote), all of which produce 2D vector graphics.
The framework has not been tested on other visual code paradigms such as Processing, p5.js, or raster-generating code, nor on 3D formats.
The applicability of our findings to these broader settings is unknown.

\paragraph{(f) Negative control scope.}
Our pre-registered C1 negative control (\S\ref{sec:c1}) formally fails the D060 criterion: non-visual Python code yields 60--110\% of the visual cross-format CKA in magnitude across 3~models $\times$ 7~layers.
The Python snippets are randomly sampled (not semantically matched to the visual triples), which limits the precision of this comparison.
A stronger control would use semantically matched non-visual code triples (e.g., Python programs describing analogous computational structures), which remains a non-trivial dataset construction challenge.
This failure establishes that the measured cross-format signal reflects code-general rather than visual-specific structure; the random sampling limits only the precision of this quantification, not the qualitative conclusion.
